\begin{thebibliography}{99}
\footnotesize\itemsep=0pt

\bibitem{allegretti2019categorified}
Allegretti D.G.L., Categorified canonical bases and framed {BPS} states,
 \href{https://doi.org/10.1007/s00029-019-0518-3}{\textit{Selecta Math.~(N.~S.)}} \textbf{25} (2019), 69, 50~pages,
 \href{https://arxiv.org/abs/1806.10394}{arXiv:1806.10394}.

\bibitem{BerensteinFominZelevinsky96}
Berenstein A., Fomin S., Zelevinsky A., Parametrizations of canonical bases and
 totally positive matrices, \href{https://doi.org/10.1006/aima.1996.0057}{\textit{Adv. Math.}} \textbf{122} (1996), 49--149.

\bibitem{BerensteinFominZelevinsky05}
Berenstein A., Fomin S., Zelevinsky A., Cluster algebras~{III}:~{U}pper bounds
 and double {B}ruhat cells, \href{https://doi.org/10.1215/S0012-7094-04-12611-9}{\textit{Duke Math.~J.}} \textbf{126} (2005), 1--52,
 \href{https://arxiv.org/abs/math.RT/0305434}{arXiv:math.RT/0305434}.

\bibitem{berenstein1997total}
Berenstein A., Zelevinsky A., Total positivity in {S}chubert varieties,
 \href{https://doi.org/10.1007/PL00000363}{\textit{Comment. Math. Helv.}} \textbf{72} (1997), 128--166.

\bibitem{BerensteinZelevinsky05}
Berenstein A., Zelevinsky A., Quantum cluster algebras, \href{https://doi.org/10.1016/j.aim.2004.08.003}{\textit{Adv. Math.}}
 \textbf{195} (2005), 405--455, \href{https://arxiv.org/abs/math.QA/0404446}{arXiv:math.QA/0404446}.

\bibitem{chang2019cluster}
Chang W., Schiffler R., Cluster automorphisms and quasi-automorphisms,
 \href{https://doi.org/10.1016/j.aam.2019.07.007}{\textit{Adv. in Appl. Math.}} \textbf{110} (2019), 342--374,
 \href{https://arxiv.org/abs/1808.02108}{arXiv:1808.02108}.

\bibitem{DerksenWeymanZelevinsky09}
Derksen H., Weyman J., Zelevinsky A., Quivers with potentials and their
 representations~{II}: {A}pplications to cluster algebras, \href{https://doi.org/10.1090/S0894-0347-10-00662-4}{\textit{J.~Amer.
 Math. Soc.}} \textbf{23} (2010), 749--790, \href{https://arxiv.org/abs/0904.0676}{arXiv:0904.0676}.

\bibitem{dupont2011generic}
Dupont G., Generic variables in acyclic cluster algebras, \href{https://doi.org/10.1016/j.jpaa.2010.06.012}{\textit{J.~Pure Appl.
 Algebra}} \textbf{215} (2011), 628--641, \href{https://arxiv.org/abs/1006.0166}{arXiv:1006.0166}.

\bibitem{FockGoncharov06a}
Fock V.V., Goncharov A.B., Moduli spaces of local systems and higher
 {T}eichm\"uller theory, \href{https://doi.org/10.1007/s10240-006-0039-4}{\textit{Publ. Math. Inst. Hautes \'Etudes Sci.}}
 (2006), 1--211, \href{https://arxiv.org/abs/math.AG/0311149}{arXiv:math.AG/0311149}.

\bibitem{FockGoncharov09}
Fock V.V., Goncharov A.B., Cluster ensembles, quantization and the dilogarithm,
 \href{https://doi.org/10.24033/asens.2112}{\textit{Ann. Sci. \'Ec. Norm. Sup\'er.}} \textbf{42} (2009), 865--930,
 \href{https://arxiv.org/abs/math.AG/0311245}{arXiv:math.AG/0311245}.

\bibitem{FockGoncharov09a}
Fock V.V., Goncharov A.B., The quantum dilogarithm and representations of
 quantum cluster varieties, \href{https://doi.org/10.1007/s00222-008-0149-3}{\textit{Invent. Math.}} \textbf{175} (2009),
 223--286, \href{https://arxiv.org/abs/math.QA/0702397}{arXiv:math.QA/0702397}.

\bibitem{FominShapiroThurston08}
Fomin S., Shapiro M., Thurston D., Cluster algebras and triangulated surfaces.
 {P}art~{I}: {C}luster complexes, \href{https://doi.org/10.1007/s11511-008-0030-7}{\textit{Acta Math.}} \textbf{201} (2008),
 83--146, \href{https://arxiv.org/abs/math.RA/0608367}{arXiv:math.RA/0608367}.

\bibitem{FominZelevinsky02}
Fomin S., Zelevinsky A., Cluster algebras~{I}: {F}oundations, \href{https://doi.org/10.1090/S0894-0347-01-00385-X}{\textit{J.~Amer.
 Math. Soc.}} \textbf{15} (2002), 497--529, \href{https://arxiv.org/abs/math.RT/0104151}{arXiv:math.RT/0104151}.

\bibitem{fraser2016quasi}
Fraser C., Quasi-homomorphisms of cluster algebras, \href{https://doi.org/10.1016/j.aam.2016.06.005}{\textit{Adv. in Appl.
 Math.}} \textbf{81} (2016), 40--77, \href{https://arxiv.org/abs/1509.05385}{arXiv:1509.05385}.

\bibitem{GeissLeclercSchroeer10}
Gei{\ss} C., Leclerc B., Schr\"oer J., Kac--{M}oody groups and cluster
 algebras, \href{https://doi.org/10.1016/j.aim.2011.05.011}{\textit{Adv. Math.}} \textbf{228} (2011), 329--433,
 \href{https://arxiv.org/abs/1001.3545}{arXiv:1001.3545}.

\bibitem{GeissLeclercSchroeer10b}
Gei{\ss} C., Leclerc B., Schr\"oer J., Generic bases for cluster algebras and
 the {C}hamber {A}nsatz, \href{https://doi.org/10.1090/S0894-0347-2011-00715-7}{\textit{J.~Amer. Math. Soc.}} \textbf{25} (2012),
 21--76, \href{https://arxiv.org/abs/1004.2781}{arXiv:1004.2781}.

\bibitem{GeissLeclercSchroeer11}
Gei{\ss} C., Leclerc B., Schr\"oer J., Cluster structures on quantum coordinate
 rings, \href{https://doi.org/10.1007/s00029-012-0099-x}{\textit{Selecta Math.~(N.~S.)}} \textbf{19} (2013), 337--397,
 \href{https://arxiv.org/abs/1104.0531}{arXiv:1104.0531}.

\bibitem{gekhtman2017hamiltonian}
Gekhtman M., Nakanishi T., Rupel D., Hamiltonian and {L}agrangian formalisms of
 mutations in cluster algebras and application to dilogarithm identities,
 \href{https://doi.org/10.1093/integr/xyx005}{\textit{J.~Integrable Syst.}} \textbf{2} (2017), xyx005, 35~pages,
 \href{https://arxiv.org/abs/1611.02813}{arXiv:1611.02813}.

\bibitem{GekhtmanShapiroVainshtein03}
Gekhtman M., Shapiro M., Vainshtein A., Cluster algebras and {P}oisson
 geometry, \href{https://doi.org/10.17323/1609-4514-2003-3-3-899-934}{\textit{Mosc. Math.~J.}} \textbf{3} (2003), 899--934,
 \href{https://arxiv.org/abs/math.QA/0208033}{arXiv:math.QA/0208033}.

\bibitem{GekhtmanShapiroVainshtein05}
Gekhtman M., Shapiro M., Vainshtein A., Cluster algebras and
 {W}eil--{P}etersson forms, \href{https://doi.org/10.1215/S0012-7094-04-12723-X}{\textit{Duke Math.~J.}} \textbf{127} (2005),
 291--311, \href{https://arxiv.org/abs/math.QA/0309138}{arXiv:math.QA/0309138}.

\bibitem{GY13}
Goodearl K.R., Yakimov M.T., Quantum cluster algebra structures on quantum
 nilpotent algebras, \href{https://doi.org/10.1090/memo/1169}{\textit{Mem. Amer. Math. Soc.}} \textbf{247} (2017),
 vii+119~pages, \href{https://arxiv.org/abs/1309.7869}{arXiv:1309.7869}.

\bibitem{goodearl2016berenstein}
Goodearl K.R., Yakimov M.T., The {B}erenstein--{Z}elevinsky quantum cluster
 algebra conjecture, \href{https://doi.org/10.4171/JEMS/969}{\textit{J.~Eur. Math. Soc. (JEMS)}} \textbf{22} (2020),
 2453--2509, \href{https://arxiv.org/abs/1602.00498}{arXiv:1602.00498}.

\bibitem{goodearl2020integral}
Goodearl K.R., Yakimov M.T., Integral quantum cluster structures, \href{https://doi.org/10.1215/00127094-2020-0061}{\textit{Duke
 Math.~J.}} \textbf{170} (2021), 1137--1200, \href{https://arxiv.org/abs/2003.04434}{arXiv:2003.04434}.

\bibitem{gross2013birational}
Gross M., Hacking P., Keel S., Birational geometry of cluster algebras,
 \href{https://doi.org/10.14231/AG-2015-007}{\textit{Algebr. Geom.}} \textbf{2} (2015), 137--175, \href{https://arxiv.org/abs/1309.2573}{arXiv:1309.2573}.

\bibitem{gross2018canonical}
Gross M., Hacking P., Keel S., Kontsevich M., Canonical bases for cluster
 algebras, \href{https://doi.org/10.1090/jams/890}{\textit{J.~Amer. Math. Soc.}} \textbf{31} (2018), 497--608,
 \href{https://arxiv.org/abs/1411.1394}{arXiv:1411.1394}.

\bibitem{ip2018cluster}
Ip I.C.H., Cluster realization of {$\mathcal{U_q(\mathfrak{g})}$} and
 factorizations of the universal {$R$}-matrix, \href{https://doi.org/10.1007/s00029-018-0432-0}{\textit{Selecta Math.~(N.~S.)}}
 \textbf{24} (2018), 4461--4553, \href{https://arxiv.org/abs/1612.05641}{arXiv:1612.05641}.

\bibitem{Kang2018}
Kang S.-J., Kashiwara M., Kim M., Oh S.-J., Monoidal categorification of cluster
 algebras, \href{https://doi.org/10.1090/jams/895}{\textit{J.~Amer. Math. Soc.}} \textbf{31} (2018), 349--426,
 \href{https://arxiv.org/abs/1412.8106}{arXiv:1412.8106}.

\bibitem{Kas:crystal}
Kashiwara M., On crystal bases of the {$Q$}-analogue of universal enveloping
 algebras, \href{https://doi.org/10.1215/S0012-7094-91-06321-0}{\textit{Duke Math.~J.}} \textbf{63} (1991), 465--516.

\bibitem{kashiwara2019laurent}
Kashiwara M., Kim M., Laurent phenomenon and simple modules of quiver {H}ecke
 algebras, \href{https://doi.org/10.1112/s0010437x19007565}{\textit{Compos. Math.}} \textbf{155} (2019), 2263--2295,
 \href{https://arxiv.org/abs/1811.02237}{arXiv:1811.02237}.

\bibitem{keller2011cluster}
Keller B., On cluster theory and quantum dilogarithm identities, in
 Representations of {A}lgebras and {R}elated {T}opics, \textit{EMS Ser. Congr. Rep.},
 \href{https://doi.org/10.4171/101-1/3}{Eur. Math. Soc.}, Z\"urich, 2011, 85--116, \href{https://arxiv.org/abs/1102.4148}{arXiv:1102.4148}.

\bibitem{Keller12}
Keller B., Cluster algebras and derived categories, in Derived {C}ategories in
 {A}lgebraic {G}eometry, \textit{EMS Ser. Congr. Rep.}, \href{https://doi.org/10.4171/115-1/6}{Eur. Math. Soc.}, Z\"urich,
 2012, 123--183, \href{https://arxiv.org/abs/1202.4161}{arXiv:1202.4161}.

\bibitem{Kimura10}
Kimura Y., Quantum unipotent subgroup and dual canonical basis,
 \href{https://doi.org/10.1215/21562261-1550976}{\textit{Kyoto~J. Math.}} \textbf{52} (2012), 277--331, \href{https://arxiv.org/abs/1010.4242}{arXiv:1010.4242}.

\bibitem{kimura2017twist}
Kimura Y., Oya H., Twist automorphisms on quantum unipotent cells and dual
 canonical bases, \href{https://doi.org/10.1093/imrn/rnz040}{\textit{Int. Math. Res. Not.}} \textbf{2021} (2021),
 6772--6847, \href{https://arxiv.org/abs/1701.02268}{arXiv:1701.02268}.

\bibitem{Lusztig90}
Lusztig G., Canonical bases arising from quantized enveloping algebras,
 \href{https://doi.org/10.2307/1990961}{\textit{J.~Amer. Math. Soc.}} \textbf{3} (1990), 447--498.

\bibitem{Lusztig91}
Lusztig G., Quivers, perverse sheaves, and quantized enveloping algebras,
 \href{https://doi.org/10.2307/2939279}{\textit{J.~Amer. Math. Soc.}} \textbf{4} (1991), 365--421.

\bibitem{lusztig2000semicanonical}
Lusztig G., Semicanonical bases arising from enveloping algebras, \href{https://doi.org/10.1006/aima.1999.1873}{\textit{Adv.
 Math.}} \textbf{151} (2000), 129--139.

\bibitem{mcnamara2021cluster}
McNamara P., Cluster monomials are dual canonical, \href{https://arxiv.org/abs/2112.04109}{arXiv:2112.04109}.

\bibitem{muller2015existence}
Muller G., The existence of a maximal green sequence is not invariant under
 quiver mutation, \href{https://doi.org/10.37236/5412}{\textit{Electron.~J. Combin.}} \textbf{23} (2016), 2.47,
 23~pages, \href{https://arxiv.org/abs/1503.04675}{arXiv:1503.04675}.

\bibitem{Muller12}
Muller G., Skein and cluster algebras of marked surfaces, \href{https://doi.org/10.4171/QT/79}{\textit{Quantum
 Topol.}} \textbf{7} (2016), 435--503, \href{https://arxiv.org/abs/1204.0020}{arXiv:1204.0020}.

\bibitem{Nagao10}
Nagao K., Donaldson--{T}homas theory and cluster algebras, \href{https://doi.org/10.1215/00127094-2142753}{\textit{Duke
 Math.~J.}} \textbf{162} (2013), 1313--1367, \href{https://arxiv.org/abs/1002.4884}{arXiv:1002.4884}.

\bibitem{nakanishi2012tropical}
Nakanishi T., Zelevinsky A., On tropical dualities in cluster algebras, in
 Algebraic {G}roups and {G}uantum {G}roups, \textit{Contemp. Math.}, Vol. 565,
 \href{https://doi.org/10.1090/conm/565/11159}{American Mathematical Society}, Providence, RI, 2012, 217--226,
 \href{https://arxiv.org/abs/1101.3736}{arXiv:1101.3736}.

\bibitem{plamondon2013generic}
Plamondon P.-G., Generic bases for cluster algebras from the cluster category,
 \href{https://doi.org/10.1093/imrn/rns102}{\textit{Int. Math. Res. Not.}} \textbf{2013} (2013), 2368--2420,
 \href{https://arxiv.org/abs/1111.4431}{arXiv:1111.4431}.

\bibitem{Qin12}
Qin F., {$t$}-analog of {$q$}-characters, bases of quantum cluster algebras,
 and a correction technique, \href{https://doi.org/10.1093/imrn/rnt115}{\textit{Int. Math. Res. Not.}} \textbf{2014}
 (2014), 6175--6232, \href{https://arxiv.org/abs/1207.6604}{arXiv:1207.6604}.

\bibitem{qin2017triangular}
Qin F., Triangular bases in quantum cluster algebras and monoidal
 categorification conjectures, \href{https://doi.org/10.1215/00127094-2017-0006}{\textit{Duke Math.~J.}} \textbf{166} (2017),
 2337--2442, \href{https://arxiv.org/abs/1501.04085}{arXiv:1501.04085}.

\bibitem{qin2019bases}
Qin F., Bases for upper cluster algebras and tropical points, \href{https://doi.org/10.4171/JEMS/1308}{\textit{J.~Eur.
 Math. Soc. (JEMS)}}, {t}o appear, \href{https://arxiv.org/abs/1902.09507}{arXiv:1902.09507}.

\bibitem{qin2021cluster}
Qin F., Cluster algebras and their bases, \href{https://arxiv.org/abs/2108.09279}{arXiv:2108.09279}.

\bibitem{qin2020dual}
Qin F., Dual canonical bases and quantum cluster algebras, \href{https://arxiv.org/abs/2003.13674}{arXiv:2003.13674}.

\bibitem{schrader2019k}
Schrader G., Shapiro A., {$K$}-theoretic Coulomb branches of quiver gauge
 theories and cluster varieties, \href{https://arxiv.org/abs/1910.03186}{arXiv:1910.03186}.

\bibitem{schrader2019cluster}
Schrader G., Shapiro A., A cluster realization of {$U_q(\mathfrak{sl}_n)$} from
 quantum character varieties, \href{https://doi.org/10.1007/s00222-019-00857-6}{\textit{Invent. Math.}} \textbf{216} (2019),
 799--846, \href{https://arxiv.org/abs/1607.00271}{arXiv:1607.00271}.

\bibitem{scott2006grassmannians}
Scott J.S., Grassmannians and cluster algebras, \href{https://doi.org/10.1112/S0024611505015571}{\textit{Proc. London Math.
 Soc.}} \textbf{92} (2006), 345--380, \href{https://arxiv.org/abs/math.CO/0311148}{arXiv:math.CO/0311148}.

\bibitem{shen2022cluster}
Shen L., Cluster nature of quantum groups, \href{https://arxiv.org/abs/2209.06258}{arXiv:2209.06258}.

\bibitem{Tran09}
Tran T., {$F$}-polynomials in quantum cluster algebras, \href{https://doi.org/10.1007/s10468-010-9226-6}{\textit{Algebr.
 Represent. Theory}} \textbf{14} (2011), 1025--1061, \href{https://arxiv.org/abs/0904.3291}{arXiv:0904.3291}.

\end{thebibliography}
